\chapter{Statement of Application}
\label{chap:soa}

This chapter maps each risk identified in the assessment to the relevant ISO/IEC~27001:2022 Annex~A controls\cite{ismsonline2025annexa} and documents whether a control is \textit{included} or \textit{excluded}.  For clarity, the items are grouped into two categories:

\begin{itemize}
    \item \textbf{Shared Assets} – Information‑system components and services that are shared across the organisation.
    \item \textbf{Departmental Assets} – Assets that belong to, or are primarily managed by, a particular department.
\end{itemize}

\section{Shared Assets}

The common assets listed in this section underpin the organisation's overall business operations - finance, people, 
customer engagement and core \ac{IT} infrastructure.  Because a single vulnerability can affect multiple functions, 
the controls are selected not only to mitigate the specific risk, but also to strengthen the resilience of the 
broader environment.  The table below shows how each identified threat is addressed, together with the rationale for 
including or legitimately excluding each Appendix A control.

\begin{longtable}{|p{2cm}|p{4cm}|p{4.5cm}|p{4cm}|}
    \caption{Statement of application for shared assets}\\
    \hline
    \textbf{Asset ID} & \textbf{Threat} & \textbf{Control(s)\,(Annex A)} & \textbf{Justification for Inclusion/Exclusion} \\
    \hline
    \endhead
    ERP-301  & Data‑breach (access‑control failure)              & A.9.1.1, A.9.4.1 & Critical system; access‑policy enforcement mandatory. \\
    \hline
    IT-201   & Phishing / social‑engineering attack              & A.6.3.1, A.7.2.2 & User awareness is primary defence; inclusion essential. \\
    \hline
    IT-202   & Network misconfiguration                          & A.13.1.1, A.12.4.1 & Network segmentation and logging indispensable. \\
    \hline
    CS-401   & Data leakage (excessive permissions)              & A.9.2.3, A.9.2.5 & Periodic permission reviews prevent leakage. \\
    \hline
    HR-501   & Insider threat (malicious privileged user)        & A.9.2.5, A.12.4.3 & Privileged‑activity oversight is critical. \\
    \hline
    HR-502   & Staff unavailability / single point of knowledge  & A.17.1.1, A.5.30  & Succession planning and redundancy required. \\
    \hline
\end{longtable}

Taken together, the controls selected strengthen the confidentiality, integrity and availability of the organisation's
 most widely used assets.  \emph{ERP-301} sits at the heart of finance, HR and supply chain workflows, so strengthening 
 access governance via A.9 is non-negotiable - a compromised account could otherwise expose the entire corporate dataset.
\emph{IT-201} illustrates that social engineering remains a quick way to bypass technical safeguards; ongoing awareness programmes
(A.6.3.1) and scheduled refresher training (A.7.2.2) directly reduce the success rate of such attacks.  For \emph{IT-202}, 
granular network zoning (A.13.1.1) combined with correlated event logging (A.12.4.1) enables rapid detection and isolation of 
any misconfigured firewall rule.  The customer records server, \emph{CS-401}, relies on strict least privilege (A.9.2.x); 
quarterly access reviews eliminate dormant high-privilege accounts that could leak personally identifiable information.  
Privileged HR roles in \emph{HR-501} handle payroll and medical records; extensive logging of admin sessions (A.12.4.3) 
provides both deterrence and evidence.  Finally, \emph{HR-502} addresses business continuity: cross-training and an up-to-date 
succession matrix (A.17.1.1, A.5.30) safeguard critical knowledge and ensure uninterrupted operations during pandemics or sudden 
staff departures.


\section{\ac{RD} Department}

\subsection{Control Mapping and Justification}

The following table maps selected Annex A controls to the \ac{RD} department's risks, as identified in the risk assessment.

\begin{tabular}{|p{5cm}|p{8cm}|p{2.5cm}|}
\hline
\textbf{Control (Annex A)} & \textbf{Justification} & \textbf{Status} \\
\hline
A.5.1.1 Policies for Information Security & \ac{RD} requires tailored information security policies to address specific asset and research data needs. & Included \\
\hline
A.5.18 Access Control Policy & Ensures RBAC is enforced across systems handling proprietary algorithms and IP. & Included \\
\hline
A.5.20 Authentication Information & \ac{MFA} is implemented for all critical systems. & Included \\
\hline
A.5.28 Protection of Information in Transit & Required to protect sensitive research data transmitted between systems and collaborators. & Included \\
\hline
A.5.32 Logging & Supports forensic investigation in the event of insider threats. & Included \\
\hline
A.5.36 Backup & Regular backups of research and simulation environments are necessary to protect against corruption or loss. & Included \\
\hline
A.5.25 Cryptographic Controls & Ensures encryption of research databases and development repositories. & Included \\
\hline
A.5.10 Acceptable Use of Information & Not directly applicable to \ac{RD} systems under strict access policies. & Excluded \\
\hline
\end{tabular}



\section{Supply-Chain Assets}

The supply chain assets focus on the organisation's \texttt{ERP-301} platform, which acts as a central procurement hub and integration point for dozens of external suppliers.  
Because a single weakness here can quickly propagate to production systems and disrupt service delivery, the associated controls are chosen both to mitigate the immediate threat and to harden the broader ecosystem.  
The table \ref{tab:supply-chain-soa} maps each risk to its corresponding Appendix A control and explains why the control is either included or excluded.

\begin{longtable}{|p{2cm}|p{4cm}|p{4.5cm}|p{4cm}|}
    \caption{Statement of application for supply-chain assets}
    \label{tab:supply-chain-soa}\\
    \hline
    \textbf{Asset ID} & \textbf{Threat} & \textbf{Control(s)\,(Annex A)} & \textbf{Justification for Inclusion/Exclusion} \\
    \hline
    \endhead
    ERP-301 & Supply-chain compromise (third-party)   & A.15.1.1, A.15.2.1 & Supplier vetting and continuous monitoring are mandatory. \\
    \hline
    ERP-301 & Data-integrity loss (bad patch / config) & A.12.1.2, A.14.2.2 & Formal change and update control essential. \\
    \hline
    ERP-301 & System downtime (availability loss)      & A.17.1.2, A.17.2.1 & Backup and fail-over required for operational continuity. \\
    \hline
\end{longtable}

The shared \texttt{ERP-301} instance also drives procurement, making it the single integration point with numerous vendors.  
A lapse in supplier diligence could embed malware in firmware deliveries, so controls A.15.1.1 and A.15.2.1 enforce binding security clauses, periodic assessments and service level triggers to detect anomalies early.  
Patch failures remain a major cause of \ac{ERP} data corruption; segregated test environments and dual approval workflows (A.12.1.2, A.14.2.2) minimise the chance of bad updates reaching production.  
Finally, every minute of \ac{ERP} downtime delays purchase orders and network maintenance spares, so hot standby infrastructure (A.17.1.2) and geographic replication (A.17.2.1) are essential to maintain supply chain resilience.